% function_spaces.tex
% Section 5 -- function spaces on the bipartite octahedral lattice.
% Discrete L^p / ell^p, Sobolev-analogue spaces ell^{k,p}, n-dependent
% scaling, embedding theorems.

\label{sec:function_spaces}

\placeholder{Section body -- write during Phase 3.  Define the
discrete $\ell^p$ and Sobolev-analogue $\ell^{k,p}$ spaces on the
lattice; state and prove $n$-dependent scaling, density, and
embedding theorems.}

\subsection{Discrete $\ell^p$ Spaces}
\label{subsec:lp_spaces}

\placeholder{Define $\ell^p(\Lambda_n)$ for $1 \leq p \leq \infty$
using the counting measure on $\Lambda_n$.  Standard
Banach-space properties: completeness, separability for $p <
\infty$, duality.}

\subsection{Sobolev-Analogue Spaces}
\label{subsec:sobolev_analogue}

\placeholder{Define $\ell^{k,p}(\Lambda_n)$ as the space of lattice
functions whose discrete-derivatives up to order $k$ (via the
discrete exterior derivative of
Section~\ref{sec:discrete_diff_geometry}) are in $\ell^p$.  State
the basic structural theorems (completeness, density of
finitely-supported functions, duality).}

\subsection{Scaling and Embedding Theorems}
\label{subsec:embeddings}

\placeholder{$n$-dependent scaling of the norms under
lattice-spacing rescaling $\epsilon$.  State discrete analogues of
the Sobolev embedding $W^{k,p} \hookrightarrow L^q$ with explicit
$n$-dependent exponent relations.  These are the embeddings that
will pass to the continuum limit in
Section~\ref{sec:continuum_limits}.}

\subsection{Compact Embeddings}
\label{subsec:compact_embeddings}

\placeholder{Compactness statements analogous to Rellich-Kondrachov
in the discrete setting; the role of the lattice-spacing parameter
$\epsilon$ in controlling compactness.}
